% Source printed pp. 65-69.
\section{Hilbert 概形几何入门}\label{src:2.5}
\begin{epigraph}
Hilbert 概形的墨菲定律: 没有哪种几何可能性可怕到不能在 Hilbert 概形的某个分支上一般地出现.
\epigraphauthor{Joe Harris 和 Ian Morrison~\cite[第 $18$ 页]{HM98}}
\end{epigraph}
Hilbert 概形是研究最充分的一类模空间; 过去 $50$ 年中, 受到更多关注的也许只有 $\overline M_g$ 和 $\cBun_{r,d}^{\mathrm{ps}}(C)$.
因此, 我们不打算系统论述, 只介绍几个有趣的例子和特点.
\subsection{初步例子}\label{src:2.5.1-subsection}
本节在代数闭域 $k$ 上工作.
射影概形 $X\subset\mathbb P^n$ 的 Hilbert 多项式 $P(z)=\sum_{i=0}^da_iz^i$ 编码了 $X$ 的不变量.
例如, $\dim X$ 是 $P$ 的次数 $d$, $\deg X$ 是经归一化的首项系数 $d!a_d$.
Riemann--Roch 定理表明: 对光滑曲线 $C\subseteq\mathbb P^n$, $P(z)=\deg(C)z+(1-g)$; 对光滑曲面 $S\subset\mathbb P^n$, 
\[
P(z)=\frac12(zH\cdot(zH-K))+(1-p_a),
\]
其中 $H$ 为超平面除子, $K$ 为典范除子, $p_a=1-\chi(\mathcal O_S)$ 为算术亏格.
在任意维数中, Hirzebruch--Riemann--Roch 定理断言
$P(z)=\int_X\operatorname{ch}(\mathcal O_X(z))\operatorname{td}(X)$, 
其中 $\operatorname{ch}(\mathcal O_X(z))$ 为 Chern 特征, $\operatorname{td}(X)$ 为 Todd 类.
\begin{exercise}[超曲面与线性子空间]\label{src:2.5.1}
\begin{enumerate}[(a)]
\item 次数为 $d$ 的超曲面 $H\subseteq\mathbb P^n$ 的 Hilbert 多项式为
\[
P(z)=\chi(\mathcal O_{\mathbb P^n}(z))-\chi(\mathcal O_{\mathbb P^n}(z-d))
=\binom{n+z}{n}-\binom{n+z-d}{n}.
\]
证明 $\Hilb^P(\mathbb P^n)\cong\mathbb P(\Gamma(\mathbb P^n,\mathcal O(d)))$.
\item 维数为 $k$ 的线性子空间 $L\subseteq\mathbb P^n$ 的 Hilbert 多项式为 $P(z)=\binom{z+k}{k}$, 且 $\Hilb^P(\mathbb P^n)=\Gr(k+1,n+1)$.
\end{enumerate}
\end{exercise}
\begin{example}[曲线上点的 Hilbert 概形]\label{src:2.5.2}
若 $C$ 为光滑不可约射影曲线, 证明 $n$ 个点的 Hilbert 概形 $\Hilb^n(C)$(将 $n$ 看作常数多项式)为光滑不可约射影簇, 并同构于对称积
\[
\operatorname{Sym}^nC:=
\underbrace{C\times\cdots\times C}_{n}/S_n,
\]
其中 $S_n$ 通过置换各因子作用.
注意, 这一商作为射影簇存在, 因为 $C\times\cdots\times C$ 射影; 见习题~\bookref{5.2.14}.
\end{example}
\begin{example}[曲面上点的 Hilbert 概形]\label{src:2.5.3}
若 $S$ 为光滑不可约射影曲面, 则 $n$ 个点的 Hilbert 概形 $\Hilb^n(S)$ 为光滑不可约射影簇~\cite{Fog68}.
另见~\cite{Nak99a}~和~\cite[\S$4$]{Mac07}.
存在射影簇的双有理态射
\[
\Hilb^n(S)\longrightarrow\operatorname{Sym}^n(S):=
\underbrace{S\times\cdots\times S}_{n}/S_n.
\]
当 $n>1$ 时, 对称积 $\operatorname{Sym}^n(S)$ 不光滑, 因此这给出一个奇点消解.
对无序的(可能不互异的)点组 $(p_1,\ldots,p_n)\in\operatorname{Sym}^n(S)$, 纤维由 $\{p_1,\ldots,p_n\}$ 上所有可能的长度为 $n$ 的概形结构组成.

$n=2$ 是第一个有趣的情形, 因为 $\Hilb^1(S)=S$.
对任意点 $p\in S$, 有许多支于 $p$ 的长度为 $2$ 的非约化概形结构.
它们由其``切方向''参数化: 给定 $[a:b]\in\mathbb P^1$, 
$k[x,y]/(x^2,xy,y^2,ay-bx)$ 定义一个长度为 $2$ 的非约化子概形.
此时, $\Hilb^2(S)\to\operatorname{Sym}^2(S)$ 是沿由 $p\mapsto(p,p)$ 给出的对角线 $S\hookrightarrow\operatorname{Sym}^2(S)$ 的爆破.
对 $n>2$, 映射 $\Hilb^n(S)\to\operatorname{Sym}^n(S)$ 是沿某个理想层的爆破~\cite{Hai98}, 但该理想层的描述更复杂.
当 $X$ 具有任意维数时, $\Hilb^n(X)$ 在由 $X$ 的 $n$ 个互异光滑点组成的(约化)闭子概形 $Z\subseteq X$ 处光滑.
若 $X$ 约化, 则 $\Hilb^n(X)$ 有一个维数为 $n\dim(X)$ 的开子概形, 参数化 $n$ 个互异光滑点.
Fogarty 的另一个结果是: 只要 $X$ 连通, $\Hilb^n(X)$ 就连通~\cite{Fog68}.
此外, 对每个射影概形 $X$, 存在一个称为``好分支''的不可约分支 $\Hilb^n(X)$, 可等同于 $\operatorname{Sym}^n(X)$ 沿某个理想层的爆破~\cite{ES14}.
\end{example}
\begin{example}[扭曲三次曲线]\label{src:2.5.4}
Hilbert 概形 $\Hilb^{3z+1}(\mathbb P^3)$ 是两个光滑有理不可约分支 $H$ 和 $H'$ 的并.
两者维数分别为 $12$ 和 $15$, 沿一个维数为 $11$ 的光滑有理子簇横截相交~\cite{PS85}.
轨迹 $H$ 为扭曲三次曲线组成的轨迹 $H_0$ 的闭包; 扭曲三次曲线是 $\mathbb P^3$ 中次数为 $3$ 的有理光滑曲线.
每条扭曲三次曲线都可由线丛 $\mathcal O_{\mathbb P^1}(3)$ 及
$\Gamma(\mathbb P^1,\mathcal O_{\mathbb P^1}(3))$ 的一组基给出的映射 $\mathbb P^1\to\mathbb P^3$ 表示, 且这种表示在相差 $\mathbb P^1$ 的自同构的意义下唯一.
所有这样的曲线射影等价, 即相差 $\mathbb P^3$ 的自同构, 故 $H_0$ 等同于齐性空间
$\operatorname{Aut}(\mathbb P^3)/\operatorname{Aut}(\mathbb P^1)
=\operatorname{PGL}_4/\operatorname{PGL}_2$, 
它光滑, 不可约, 维数为 $12$.
轨迹 $H_0$ 不固有, 因为它包含诸如由
$[x,y]\mapsto[x^3,x^2y,xy^2,ty^3]$
给出的 $\mathbb P^1\hookrightarrow\mathbb P^3$ 的族, 以 $t\in\mathbb A^1$ 参数化.
其极限为奇异曲线 $C_0$, 支于 $V(w)=\mathbb P^2$ 中的一条结点三次曲线($w$ 为第 $4$ 个坐标), 但在结点处带有一个嵌入点; 见~\cite[例 $9.8.4$]{Har77}.
\begin{figure}[H]
\centering\begin{tikzpicture}[x=1cm,y=1cm]
\draw[thick] (-.9,-2.3) .. controls (-.5,-1.1) and (-.65,.1) .. (-.75,1.3)
.. controls (-.15,1.8) and (.55,2.1) .. (1.05,2.15)
.. controls (1.3,.8) and (1.15,-.4) .. (.85,-1.6)
.. controls (.1,-1.7) and (-.45,-2) .. (-.9,-2.3);
\draw[thick] (-2,-.95) .. controls (-1,.05) and (-.2,.5) .. (.3,.5)
.. controls (1.4,.4) and (2.1,.4) .. (3,.5)
.. controls (1.7,.15) and (1.25,-.35) .. (.95,-.95)
.. controls (-.2,-.5) and (-1,-.65) .. (-2,-.95);
\draw[red,densely dashed] (-.7,-.68)--(1.1,.45);
\node[anchor=west,font=\small] at (1.12,2.16) {$H$(维数 $12$)};
\node[anchor=west,font=\small] at (1.6,-.5) {$H'$(维数 $15$)};
\node[red,anchor=west,font=\small] (inter) at (1.65,.95)
{$H\cap H'$(维数 $11$)};
\draw[red,->] (inter.west)--(1.15,.5);
\foreach \a/\b in {-.35/1.2,.2/1.55,.65/1.7,-.3/.8,.3/1.05,.75/1.3,
-.25/.3,.3/.55,.7/.8,-.45/-1.7,.1/-1.45,.5/-1.2,
-.35/-1.3,.25/-1.05,.65/-.85} {
\draw[blue!65,thin] (\a,\b) .. controls (\a+.12,\b+.18) and (\a-.12,\b+.05)
.. (\a+.06,\b+.23);
}
\foreach \a/\b in {-1.35/-.6,-.85/-.4,-.25/-.25,.4/-.05,1.45/.05,1.9/.22} {
\draw[blue!65,thin] (\a,\b) .. controls (\a-.16,\b+.08) and (\a+.15,\b+.1)
.. (\a,\b+.16);
\draw[blue!65,thin] (\a+.16,\b+.16) circle (.012);
}
\end{tikzpicture}

\caption{$\Hilb^{3z+1}(\mathbb P^3)$.}\label{src:2.5.5}
\end{figure}
轨迹 $H'$ 为轨迹 $H'_0$ 的闭包, 后者由子概形 $C\amalg\{p\}$ 组成, 其中 $C$ 为包含在超平面 $H$ 中的光滑三次曲线, $p\in\mathbb P^3\setminus C$.
为计算维数, 注意超平面 $H\in\mathbb P(H^0(\mathbb P^3,\mathcal O(1)))$ 的选择由 $3$ 个参数给出, 平面三次曲线 $C\in\mathbb P(H^0(H,\mathcal O_H(3)))$ 的选择由 $9$ 个参数给出, 点 $p\in\mathbb P^3\setminus C$ 的选择由 $3$ 个参数给出.
因此轨迹 $H'_0$ 光滑, 不可约, 维数为 $15$.
轨迹 $H'_0$ 不固有, 其闭包包含点 $p$ 落在曲线上时的退化.

交集 $H\cap H'$ 由在奇点处带有嵌入点的平面奇异三次曲线组成.
这一轨迹包含上述 $C_0$ 这样的曲线, 也包含更严重的退化曲线, 例如带有嵌入点的三重直线.
事实上, 每条曲线 $C\in H\cap H'$ 都射影等价于
$V(xz,yz,z^2,q(x,y,w))$ 定义的曲线, 其中 $q(x,y,w)$ 为在 $(0,0,1)$ 处有奇点的齐次三次多项式.
这依赖于 $11$ 个参数.
\end{example}
\subsection{几何性质}\label{src:2.5.2-subsection}
\begin{exercise}[局部性质]\label{src:2.5.6}
设 $X$ 为域 $k$ 上的射影概形, $F$ 为 $X$ 上的凝聚层.
\begin{enumerate}[(a)]
\item 设 $p\in\Quot^P(F/X)$ 为对应于商 $Q=F/K$ 的点.
证明 $T_p\Quot^P(F/X)\cong\operatorname{Hom}_{\mathcal O_X}(K,Q)$.
这推广了计算 Grassmann 概形切空间的习题~\bookref{2.2.10}.
\item 推出: 若 $p\in\Hilb^P(X)$ 为对应于理想层 $I$ 定义的闭子概形 $Z\subseteq X$ 的点, 则
$T_p\Hilb^P(X)\cong H^0(Z,N_{Z/X})$, 
其中 $N_{Z/X}$ 为法层 $\operatorname{Hom}_{\mathcal O_Z}(I_Z/I_Z^2,\mathcal O_Z)$.
(这恢复了命题~\bookref{C.1.3}.)
\end{enumerate}
\end{exercise}
\textit{非空性.} Hilbert 概形 $\Hilb^P(\mathbb P^n)$ 非空, 当且仅当 Hilbert 多项式 $P$ 可写成
\begin{equation}\label{src:2.5.7}
P(z)=\binom{z+\lambda_1-1}{\lambda_1-1}
+\binom{z+\lambda_2-2}{\lambda_2-1}
+\cdots+\binom{z+\lambda_r-r}{\lambda_r-1},
\end{equation}
其中 $\lambda_1\geq\lambda_2\geq\cdots\geq\lambda_r\geq1$ 为整数.
这是 Hartshorne 的结果~\cite[推论 $5.7$]{Har66b}.

\textit{连通性.} Hartshorne 连通性定理断言: 对每个 Hilbert 多项式 $P$, Hilbert 概形 $\Hilb^P(\mathbb P^n)$ 都连通~\cite{Har66b}.
证明策略是给出从任意闭子概形 $Z\subseteq\mathbb P^n$ 到由单项式理想定义的子概形 $V(I)\subseteq\mathbb P^n$ 的退化.
这将定理归约为一个组合问题: 通过 $\mathbb A^1$ 上的族连接任意两个单项式理想.
最终, 这成了纯粹的形变与组合问题; 用 Hartshorne 的话说: ``看来在证明中实际上从未需要 Hilbert 概形.''

\textit{墨菲定律.} Mumford 给出了第一个病态例子~\cite{Mum62}: $\Hilb^{14z-23}(\mathbb P^3)$ 有一个一般地非约化的不可约分支.
Ellia, Hirschowitz 和 Mezzetti 证明, $\Hilb^{az+b}(\mathbb P^3)$ 的不可约分支数不能由 $a,b$ 的多项式控制~\cite{EHM92}.
Vakil 将墨菲定律精确化~\cite{Vak06}: 
对每个 $\mathbb Z$ 上的有限型概形 $X$ 及点 $x\in X$, 某个 Hilbert 概形中存在点
$[Z\subseteq\mathbb P^n]\in\Hilb^P(\mathbb P^n)$, 
使 $(X,x)$ 与 $(\Hilb^P(\mathbb P^n),[Z\subseteq\mathbb P^n])$ 光滑局部同构, 即它们的完备局部环在添加形式幂级数环后同构.
事实上, 可以取 $[Z\subseteq\mathbb P^n]$ 为光滑曲线!
许多其他模空间也满足墨菲定律: Kontsevich 的映射模空间, 典范极化光滑曲面的模, 带线性系的曲线的模, 以及稳定层的模空间.

\textit{光滑性.} 尽管存在墨菲定律, Hilbert 概形却出人意料地经常光滑.
此前我们已看到, 光滑曲面上点的 Hilbert 概形光滑.
Skjelnes 和 Smith 的近期定理~\cite{SS20}~利用如~\eqref{src:2.5.7}~那样定义 $P$ 的分拆 $\lambda=(\lambda_1,\ldots,\lambda_r)$, 给出了 Hilbert 概形 $\Hilb^P(\mathbb P^n)$ 光滑的充要条件.
